\scanpage{17}
\mspage{8}

Si $\mathcal D$ est un groupe isomorphe à $\mathbb D_4$, l'ens $S$ des éléments d'ordre $2$ \deleted{\unsure{non centraux}} \inserted{$\ne\underline a$ (où} $\underline a\longrightarrow$ l'unique élément central $\ne1$) \deleted{est} a quatre éléments, \deleted{qui forment deux classes par translation} \inserted{sur lequel on a une involution} \inserted{sans} point fixe par
\[
x\longmapsto x+\underline a\overset{\mathrm{def}}=x'
\]
donc à $\mathcal D$ correspond un carré $Q_{\mathcal D}$, dont l'ens des sommets est $S$, et l'\deleted{\unsure{automorphisme de|automorphisme des}} \inserted{action de} l'antipodisme sur $S$ est l'involution précédente. Mais le foncteur de groupoïdes $\mathcal D\mathrel{\text{---}}Q_{\mathcal D}$
\diagram{\begin{tikzcd}[ampersand replacement=\&]
(\text{groupes isom. à }\mathbb D_4)\arrow[r]\&(\text{carrés combinatoires})
\end{tikzcd}}
est une équivalence de groupoïdes, car on reconstitue à partir de $(S,\underline a_S:\inserted{x\mapsto x'})$ $\mathcal D$ par \deleted{\unsure{purement}} la formule
\[
\mathcal D=\bigl\{\underbrace{x\in S,\ \underline a}_{\text{générateurs}}\ \big|\ 
\underbrace{[x,y]}_{\substack{\text{commutateur}\\xyx^{-1}y^{-1}}}=\underline a\ \text{ si }\ y\ne x,x'\bigr\}
\]

% slanted
\margin{\uline{NB} Si à un carré \inserted{$Q$} on associe le groupe $\mathcal D=\operatorname{Aut}Q$, puis le carré associé à ce groupe, on trouve le carré dont l'ens des sommets est $C\amalg D$, l'action de l'antipodisme dessus \deleted{\unsure{étant}} \unsure{comme orbites} $C$ et $D$, i.e. les diagonales sont $C$ et $D$. On trouve donc ``un carré avec diagonale \deleted{\unsure{privil}} distinguée\ldots''}

Si on part, pour décrire la structure de carré, de l'ens des sommets, $S$, cette structure est définie par la donnée d'un ens. \inserted{$S$} à $4$ éléments -- i.e. d'un torseur sous un vectoriel $V_S$ de rang $2$ sur $\mathbb F_2$ -- et d'une \deleted{\unsure{tr}} involution ss pts fixes de $S$, i.e. d'un $\underline a_S\in V_S^*$ \deleted{Mais se} \deleted{\unsure{$V=V_S^*$|$V_S=V_S^*$}} \inserted{(i.e. une} partition de type $(2,2)$ de $S$). Mais pour

% archive sheet number: 16

